\section{T-Chain Architecture}

\label{sec:architecture}

\subsection{ThresholdVM Overview}

\ThresholdVM{} is a native chain on Zoo Network, registered at chain creation alongside the EVM L2, Identity Chain, and Experience Ledger. It exposes a minimal state machine with three transaction types:

\begin{enumerate}
\item \textbf{KeygenTx}: Initiates a distributed key generation session among a quorum of validators.
\item \textbf{SignTx}: Requests a threshold signature over an arbitrary message digest.
\item \textbf{ReshareTx}: Re-distributes key shares to a new validator set without changing the public key.
\end{enumerate}

Each transaction type maps to a multi-round interactive protocol. \ThresholdVM{} manages session state, enforces round deadlines, and commits the protocol transcript to the chain once all rounds complete or a timeout triggers abort.

\subsection{RPC API}

The \TChain{} RPC surface is accessible at \texttt{/v1/threshold}:

\begin{table}[h]
\centering
\begin{tabular}{lp{7cm}}
\toprule
\textbf{Method} & \textbf{Description} \\
\midrule
\texttt{threshold\_keygen} & Start a $t$-of-$n$ DKG session; returns \texttt{session\_id} and public key on completion \\
\texttt{threshold\_sign} & Submit a message hash for threshold signing; returns partial signatures, then aggregated signature \\
\texttt{threshold\_reshare} & Initiate share redistribution to a new committee \\
\texttt{threshold\_status} & Query session state (pending, round $k$, complete, aborted) \\
\texttt{threshold\_pubkey} & Retrieve the group public key for a given \texttt{key\_id} \\
\bottomrule
\end{tabular}
\caption{T-Chain RPC methods}
\label{tab:rpc}
\end{table}

\subsection{Session Management}

A session $\mathcal{S} = (\texttt{session\_id}, \texttt{protocol}, t, n, \texttt{participants}, \texttt{round}, \texttt{deadline})$ tracks the lifecycle of each interactive protocol instance. Sessions are created by a \texttt{KeygenTx} or \texttt{SignTx} and progress through rounds as validators submit messages. If a validator fails to submit a round message before the deadline, the session enters \emph{identifiable abort}: the non-responding validator is recorded on-chain and subject to slashing.

\begin{definition}[Session Liveness]
A session $\mathcal{S}$ completes if at least $t$ of the $n$ participants submit valid round messages within every deadline. Formally, for each round $r \in \{1, \ldots, R\}$:
\[
|\{i \in \mathcal{S}.\texttt{participants} : \texttt{msg}_{i,r} \text{ valid} \}| \geq t
\]
\end{definition}

\subsection{Permissioned Key Creation}

Not every actor may initiate a keygen. \TChain{} enforces an \texttt{AuthorizedChains} registry: only chains listed in this registry (e.g., the EVM L2 bridge contract, the Identity Chain) may request key creation. This prevents spam and ensures that every generated key serves a concrete cross-chain or custodial purpose.

%% --------------------------------------------------------------------------
